\section{Evaluation}
\label{sec:experiments}

\subsection{Dataset: Snowy Scenes}
\label{sec:dataset}

We evaluate on Snowy Scenes~\cite{SnowyScenes2026}, a multimodal driving dataset captured entirely in real winter conditions and distributed as a single $\sim$49\,GB archive (ROADVIEW5k), which we read lazily via \texttt{zipfile} rather than extracting to disk given its $\sim$93\,GB uncompressed footprint. It contains 5{,}027 frames across three splits (3{,}016 train / 1{,}005 val / 1{,}006 test), matching the dataset's own stated frame count, with per-frame RGB camera images, 128-beam LiDAR point clouds in the same $(N,4)$ float32 layout as KITTI, 3D object boxes already expressed in the LiDAR frame (no camera$\leftrightarrow$LiDAR extrinsic transform required, unlike KITTI), and per-point semantic segmentation labels across 29 SemanticKITTI-style classes -- the segmentation labels are not consumed here, since CRAF-X's head has no segmentation output.

Two things about the real archive turned out different from what we had planned around before access arrived. First, Snowy Scenes' three weather categories -- \texttt{accumulated}, \texttt{highway}, and \texttt{falling} -- are all already-snowy driving; there is no dry/clear split to calibrate on or to serve as the nominal side of an onset transition, which the original plan had assumed would exist. Second, within-category severity does not ramp monotonically: sampling the ground-truth per-point ``snow'' semantic class fraction across a timestamp-ordered \texttt{falling} sequence shows real but noisy fluctuation ($0.002$--$0.044$) rather than a clean onset curve, since active snowfall is bursty rather than steadily accumulating. What the data does support is a real, physically grounded severity contrast \emph{across} categories: measuring the mean per-point ``snow''-class fraction per category gives $\texttt{accumulated} = 0.0001 < \texttt{highway} = 0.0041 < \texttt{falling} = 0.0119$, an ordering with a clear physical explanation -- \texttt{falling} captures active, airborne snowfall, which LiDAR registers as spurious near-range returns, while \texttt{accumulated} is settled snow on surfaces rather than particles in the air.

\subsection{Constructing a Real Onset Stream Without a Clear-Weather Split}
\label{sec:onset-stream}

We use this ordering to build \texttt{RealSnowOnsetStream}: a nominal (held-out \texttt{accumulated}) prefix spliced with a degraded (\texttt{falling}) suffix at a known onset index, drawing frames from each category's frame pool with wraparound if a requested scene length exceeds the pool size. Every individual frame on both sides of the transition is real sensor data -- real images, real LiDAR, real object labels; only the splice point itself is a construction, since Snowy Scenes' own sequences do not contain a within-sequence weather transition. We are explicit that this is a scoping compromise relative to the originally planned design (an in-dataset clear$\rightarrow$snow transition), not the design we would have chosen if the data supported it, and we flag the alternative we did not have the means to try: splicing in real KITTI clear-weather frames as the nominal side instead of held-out \texttt{accumulated} frames, which would give a genuine cross-dataset clear-to-snow contrast at the cost of mixing class spaces -- an option for future work if reviewers judge the within-Snowy-Scenes splice an insufficiently strong stand-in for a naturally occurring onset.

The same absence of a true clear-weather condition affects the false-alarm-rate control in our operating-curve evaluation (Section~\ref{sec:results}): with no clear split available, our ``clear'' replicates use \texttt{accumulated} frames on \emph{both} sides of the splice, so the stream is stationary throughout rather than transitioning at all. This measures the false-alarm rate under a real but comparatively mild, non-transitioning snow condition, not under true clear weather, and we report it as exactly that rather than implying a stronger guarantee than the data supports.

\subsection{Detector Training and a Real Negative Result}
\label{sec:training}

We train CRAF-X for 5 epochs on the real \texttt{train} split (3{,}016 frames) under its Adversarial Consistency Training (ACT) objective, on a single local NVIDIA RTX 5000 Ada GPU, at $128\times128$ BEV resolution, batch size 4, Adam with learning rate $10^{-4}$; training loss converges from $\approx 2.3$ to $\approx 0.008$. Getting a real, trained checkpoint onto real data surfaced a finding that changed our own evaluation before we could report the operating curve honestly.

Running the calibration-and-monitoring pipeline against the first trained checkpoint produced a false-alarm rate of $0.0$ across every tested $\delta$ and detection delays that were \emph{numerically identical} between the covariate-informed and covariate-blind bettors, to every digit the operating curve reports. That is not the outcome our method predicts if the covariate carries any signal at all, and treating it as a clean result rather than investigating it would have been a mistake. Sampling the CCP consistency score $S$ directly across a real onset stream with the trained model showed why: $S$ had collapsed to $\approx 0.9993$ uniformly, with standard deviation $\approx 0.0001$--$0.0002$, statistically indistinguishable between the nominal (\texttt{accumulated}) and degraded (\texttt{falling}) regimes. With $1 - S$ this close to a constant, $\text{CCPInformedBettor}$'s $(1 + \kappa \cdot \overline{(1-S)})$ scaling term barely perturbs the base bettor's own $\lambda_t$ -- the two bettors were not tied by coincidence; they were, numerically, almost the same bettor.

Tracing the cause led to \texttt{craf\_x/training/adversarial.py} rather than to anything in the monitoring code itself. The CCP contrastive loss, $L_{\mathrm{ccp}} = -\sum_{(i,j) \in M} \log S(i,j) - \sum_{(i,j) \notin M} \log(1 - S(i,j))$, was being computed exactly once per training step, on the \emph{clean} forward pass, against the dataset's mask $m$ -- which is all-ones, because real, synchronized Snowy Scenes frames have no annotated camera$\leftrightarrow$LiDAR mismatch to supervise against. The Adversarial Consistency Training loop does compute two additional CCP forward passes, on camera- and LiDAR-adversarially-perturbed features respectively, but those scores were only ever consumed by the KL-divergence attribution-consistency term, never by the CCP contrastive loss itself with a mismatched target. Nothing in the training objective, as implemented, ever penalized $S$ for staying high under a genuine cross-modal perturbation; minimizing $L_{\mathrm{ccp}}$ against an all-ones mask is trivially satisfied by pushing $S \to 1$ everywhere, and that is exactly what five epochs of training converged to.

We consider this worth reporting as a methodological finding in its own right, independent of whether the reader cares about CRAF-X specifically: a signal trained under an \emph{invariance} objective (here, adversarial \emph{consistency} training, whose whole point is to make the cross-modal agreement score robust to attack) is, by construction, in tension with reusing that same signal as a \emph{leading indicator of degradation} for a downstream monitor, which requires exactly the sensitivity the training objective was suppressing. Any paper proposing to repurpose an internal robustness or consistency signal as an external covariate for a statistically distinct purpose should check, empirically, that the signal still varies at all under the condition it is meant to detect -- we did not check this before the first real run, and the identical-bettors result is what caught it.

The fix we applied supervises the two adversarial branches with the same CCP loss, against a mismatched ($m=0$) target, alongside the existing clean-branch call: $L_{\mathrm{ccp}} = L_{\mathrm{ccp}}^{\mathrm{clean}} + L_{\mathrm{ccp}}^{\mathrm{adv\text{-}cam}} + L_{\mathrm{ccp}}^{\mathrm{adv\text{-}lid}}$, since a PGD-perturbed feature pair is a real cross-modal mismatch by construction and was already being computed for another purpose. We added a direct unit test confirming \texttt{compute\_ccp\_loss} penalizes a confidently-agreeing score under a mismatched mask more than a confidently-disagreeing one, and an integration test confirming \texttt{act\_training\_step}'s loss decomposition actually includes a nonzero adversarial-CCP term, as a regression guard against this exact collapse recurring silently. We then retrained CRAF-X from scratch under the corrected objective and re-ran the full evaluation; Section~\ref{sec:results} reports what changed.

\subsection{Evaluation Protocol}
\label{sec:protocol}

Calibration uses 40 held-out \texttt{accumulated} frames from the \texttt{val} split, disjoint from the frames used as the nominal portion of any monitored stream, at target miscoverage $\alpha = 0.2$. Each onset stream has scene length 20 with the true onset at frame 8 (nominal frames $t < 8$, degraded frames $t \ge 8$); we sweep the false-alarm budget $\delta \in \{0.3, 0.1, 0.05\}$ and report, at each $\delta$, the empirical false-alarm rate (fraction of ``clear'' replicates that alarm at all) and the mean detection delay (alarm time minus true onset frame, among replicates that alarm; the count that never alarm within the 20-frame window is reported separately as censored rather than silently dropped). We use 5 onset replicates and 5 ``clear'' replicates per condition, and the covariate-informed bettor's gain is fixed at $\kappa = 2.0$ -- both replicate count and $\kappa$ are small-scale defaults inherited from our earlier mock-data smoke testing rather than values tuned against real data, and we say so explicitly in Section~\ref{sec:conclusion} rather than presenting them as optimized. We are equally direct about what this replicate count can and cannot support statistically: $n=5$ is a pilot-scale demonstration, not a powered comparison, and we do not report a significance test between bettors' detection-delay distributions in Section~\ref{sec:results} -- an observed gap at this $n$ should be read as suggestive, not decisive. One property of our evaluation does partially mitigate this: because each replicate's wealth trajectory is computed once and is delta-independent (Section~\ref{sec:eprocess}), every $\delta$ in the sweep and both bettors are evaluated against the identical five sampled onset streams rather than independently resampled ones, so the comparison across bettors at fixed $\delta$ is a paired comparison on shared random draws, not two independent samples -- a real if partial source of variance reduction, though not a substitute for a larger $n$. The calibration set size ($n=40$) is likewise small enough that the finite-sample conformal guarantee, while still formally valid at any $n$, is loose in practice at $\alpha=0.2$; we do not correct for this and report $\hat{q}_\alpha$ as computed, rather than presenting it as a tight threshold.

\subsection{Operating Curve on Real Data}
\label{sec:results}

Table~\ref{tab:operating-curve} reports the operating curve against the corrected checkpoint (Section~\ref{sec:training}), evaluated on held-out \texttt{val}-split frames the detector never saw during training. We report this exactly as obtained; the result is not the one our method's stated rationale predicts, and Section~\ref{sec:discussion} explains why in detail rather than qualifying it away.

\begin{table}[t]
\centering
\caption{Operating curve on the real Snowy Scenes onset stream (accumulated $\to$ falling), 5 onset + 5 stationary-accumulated replicates per condition, corrected checkpoint (\texttt{checkpoint\_final.pth}, epoch 4). FAR = false-alarm rate; delay is in frames relative to the true onset at $t=8$; censored = replicates that never alarmed within the 20-frame window.}
\label{tab:operating-curve}
\begin{threeparttable}
\begin{tabular}{@{}lccccc@{}}
\toprule
Bettor & $\delta$ & FAR & Mean delay & Censored \\
\midrule
Covariate-blind (aGRAPA) & 0.30 & 0.00 & 3.0 & 0/5 \\
Covariate-blind (aGRAPA) & 0.10 & 0.00 & 10.0 & 0/5 \\
Covariate-blind (aGRAPA) & 0.05 & 0.00 & -- & 5/5 \\
CCP-informed & 0.30 & 0.00 & -- & 5/5 \\
CCP-informed & 0.10 & 0.00 & -- & 5/5 \\
CCP-informed & 0.05 & 0.00 & -- & 5/5 \\
\bottomrule
\end{tabular}
\begin{tablenotes}
\small
\item[] $q_{\hat\alpha} = 10.19$ (40 calibration frames, $\alpha=0.2$). Covariate-blind numbers are numerically unchanged from the pre-fix checkpoint (the bettor never consumes the CCP score, so this is expected, not evidence the retrain had no effect elsewhere). CCP-informed alarms on \emph{zero} of 15 onset-replicate/$\delta$ combinations after the fix, versus alarming at every $\delta$ before it (Section~\ref{sec:training}) -- the fix changed the outcome substantially, in the opposite direction from the one motivating this paper.
\end{tablenotes}
\end{threeparttable}
\end{table}

\subsection{Discussion: the Fix Did Not Help, and Why}
\label{sec:discussion}

Fixing the CCP-collapse bug did what it was supposed to do at the level of the covariate itself: sampling $S$ across the same real onset stream with the corrected checkpoint shows real, non-degenerate variation ($S \approx 0.40$, std $\approx 0.08$, in the nominal regime; $S \approx 0.44$, std $\approx 0.03$, in the degraded regime) in place of the pre-fix collapse to $S \approx 0.9993$ uniformly. The covariate is no longer constant. It is, however, informative in the wrong direction for this task: mean disagreement $\overline{(1-S)}$ is \emph{higher} in the nominal (\texttt{accumulated}) regime ($\approx 0.599$ at $t=0$, declining slightly to $\approx 0.585$ by $t=7$) than in the degraded (\texttt{falling}) regime ($\approx 0.564$, roughly flat from $t=8$ onward) -- the opposite ordering from what a leading indicator of degradation would need. Since $\text{CCPInformedBettor}$'s update, $\lambda_t = \lambda_t^{\text{base}} \cdot (1 + \kappa \cdot \overline{(1-S_t)})$, can only scale the base bettor's stake \emph{up} (Section~\ref{sec:informed}), and disagreement here sits around $0.56$--$0.60$ throughout \emph{both} regimes rather than rising specifically at onset, the covariate-informed bettor bets aggressively for the entire stream, calibration set included. A wealth process betting aggressively on frames where the true miscoverage rate $m(t)$ happens to sit below $\alpha$ loses wealth faster than a conservatively betting one does (the factor $1 + \lambda(m(t)-\alpha)$ shrinks toward zero as $\lambda$ grows whenever $m(t) < \alpha$), and that appears to be exactly what happens here: the covariate-informed process never accumulates enough wealth within the 20-frame window to cross even the most permissive alarm threshold we tested ($\delta = 0.3$), across all 5 onset replicates.

We think the most defensible explanation is a mismatch between what the fix taught the CCP score and what this paper needs it to detect. The fix supervises $S$ using PGD-adversarially-perturbed features as the mismatched-target examples (Section~\ref{sec:training}) -- a natural and, for CRAF-X's own adversarial-robustness objective, an appropriate choice, since PGD perturbations are a real cross-modal mismatch by construction. But an adversarial perturbation crafted to maximize the detector's box-regression loss under an $\ell_\infty$ budget is a qualitatively different kind of feature-space disturbance than the one real snowfall actually produces on LiDAR returns and camera contrast. Teaching $S$ to recognize the former does not obviously transfer to recognizing the latter, and our post-fix measurement suggests it did not: whatever pattern the corrected CCP head is now picking up on correlates with something in the \texttt{accumulated}-vs-\texttt{falling} contrast, but not with which side of that contrast is actually more degraded from the detector's perspective. This is a substantive negative result about the specific mechanism this paper proposed, not a bug in the monitoring code -- the identical experiment, run twice (Section~\ref{sec:training} and here), surfaced two different real problems in two different places (the training objective, then the covariate's transfer properties), and we report both rather than stopping at the first one that looked fixed.

We are not able, within this evaluation, to distinguish between two explanations for that transfer failure: that adversarial-perturbation-based CCP supervision is fundamentally the wrong kind of covariate for a naturally-occurring weather shift regardless of how it is tuned, or that $\kappa=2.0$ and the untuned architecture we trained are simply the wrong operating point and a different $\kappa$, more training data, or a covariate computed differently (e.g., a discrepancy measure between the two modalities' own predicted object locations, rather than a jointly-trained consistency head) would recover an advantage. Section~\ref{sec:conclusion} treats this as the open problem it is, rather than resolving it by assertion.
